% content/talk.tex

\section{Layouts + slots}

\begin{frame}[layout=fullA]{Full layout (main region)}
  \BTKContent{text}{
    \begin{itemize}
      \item This slide uses \texttt{layout=full} and fills the \texttt{text} textbox.
      \begin{itemize}
      \item More test text
      \end{itemize}
      \item Even spacing helpers: \texttt{\string\hfill}, \texttt{\string\vfill}, or BTK lists.
      \item Use \texttt{BTKOverlay} for absolute/TikZ power moves.
    \end{itemize}
  }

  \BTKContent{Overlay}{
    % “power move” example: draw a box, use coordinates already defined!
    \draw[color=black,fill=red] ($(NCSA@t1@nw |- NCSA@t1@se)+(1cm,3cm)$) rectangle ($(NCSA@t1@se)+(-1cm,0cm)$);
  }
\end{frame}

\begin{frame}[layout=fullB]{Full layout (main region)}
  \BTKContent{text}{
    \begin{itemize}
      \item This slide uses \texttt{layout=full} and fills the \texttt{text} textbox.
      \begin{itemize}
      \item More test text
      \end{itemize}
      \item Even spacing helpers: \texttt{\string\hfill}, \texttt{\string\vfill}, or BTK lists.
      \item Use \texttt{BTKOverlay} for absolute/TikZ power moves.
    \end{itemize}
  }

  \BTKContent{Overlay}{
    % “power move” example: draw a box, use coordinates already defined!
    \draw[color=black,fill=red] ($(NCSA@t1@nw |- NCSA@t1@se)+(1cm,3cm)$) rectangle ($(NCSA@t1@se)+(-1cm,0cm)$);
  }
\end{frame}


\begin{frame}[layout=fullA]{Full layout with plain text}
  \BTKContent{text}{
    \lipsum[1-3]
  }
\end{frame}

\begin{frame}[layout=fullB]{Full layout with plain text}
  \BTKContent{text}{
    \lipsum[1-3]
  }
\end{frame}

\begin{frame}[layout=two-colA]{Two columns (left/right regions)}
  \BTKContent{left}{
    \textbf{Left}\\[0.5em]
    \begin{itemize}
      \item bullets
      \item more bullets
    \end{itemize}
  }

  \BTKContent{right}{
    \textbf{Right}\\[0.5em]
    \begin{itemize}
      \item figures/tables fit well
      \item or equations
    \end{itemize}
  }
\end{frame}

\begin{frame}[layout=two-colB]{Two columns (left/right regions)}
  \BTKContent{left}{
    \textbf{Left}\\[0.5em]
    \begin{itemize}
      \item bullets
      \item more bullets
    \end{itemize}
  }

  \BTKContent{right}{
    \textbf{Right}\\[0.5em]
    \begin{itemize}
      \item figures/tables fit well
      \item or equations
    \end{itemize}
  }
\end{frame}

\begin{frame}[layout=img-rightA]{Text + image (img-right)}
  \BTKOption{img-width}{0.4\linewidth}
  \BTKContent{img}{
    \includegraphics[width=0.38\linewidth]{example_detail.png}
    \caption{This is a test caption}
  }
  \BTKContent{text}{
    \textbf{Detector stack}\\[0.5em]
    \begin{itemize}
      \item Silicon tracker: measures charged-particle tracks (curvature $\rightarrow p_T$)
      \item Calorimeters: measure energy deposition (EM + hadronic)
    \end{itemize}
  }
\end{frame}

\begin{frame}[layout=img-rightB]{Text + image (img-right)}
  \BTKOption{img-width}{0.4\linewidth}
  \BTKContent{img}{
    \includegraphics[width=0.38\linewidth]{example_detail.png}
    \caption{This is a test caption}
  }
  \BTKContent{text}{
    \textbf{Detector stack}\\[0.5em]
    \begin{itemize}
      \item Silicon tracker: measures charged-particle tracks (curvature $\rightarrow p_T$)
      \item Calorimeters: measure energy deposition (EM + hadronic)
    \end{itemize}
  }
\end{frame}

\begin{frame}[layout=img-rightA,bgopacity=0.35]{Custom per-frame opacity}
  \BTKOption{img-width}{0.4\linewidth}
  \BTKContent{img}{
    \includegraphics[width=0.38\linewidth]{example_detail.png}
    \caption{This is a test caption}
  }
  \BTKContent{text}{
    \textbf{Same layout, different background opacity}\\
    You can do \texttt{bg=...}, \texttt{bgfile=...}, or \texttt{bgcolor=...}.
  }
\end{frame}

\begin{frame}[layout=img-rightB,bgopacity=0.35]{Custom per-frame opacity}
  \BTKOption{img-width}{0.4\linewidth}
  \BTKContent{img}{
    \includegraphics[width=0.38\linewidth]{example_detail.png}
    \caption{This is a test caption}
  }
  \BTKContent{text}{
    \textbf{Same layout, different background opacity}\\
    You can do \texttt{bg=...}, \texttt{bgfile=...}, or \texttt{bgcolor=...}.
  }
\end{frame}

\begin{frame}[layout=blank,bgnone]{Blank master + overlay}
  % Whole page is yours. Demonstrate exposed anchors.
  \BTKContent{Overlay}{
    % Draw body rectangle
    \draw[very thick] (current page.north west) rectangle (current page.south east);

    % Place something relative to a known anchor
    \node[anchor=north west] at (current page.north west) {\Large Blank canvas};

    % Place something in the center of the frame
    \node[anchor=center] at (current page.center) {Test text in the center of the page!};

    % If a layout defined a region, you can reference it too; in blank we define main = full page
    \draw[thick] ($(current page.north west)+(1cm,-1cm)$) rectangle ($(current page.south east)+(-1cm,1cm)$);
  }
\end{frame}

\section{Conclusion}

\begin{frame}[layout=fullA]{Summary}
  \BTKContent{text}{
    \begin{itemize}
      \item Layouts define named regions.
      \item Regions give normal LaTeX flow inside bounded boxes.
      \item Overlays reference stable anchors for “Power mode”.
    \end{itemize}
  }
\end{frame}

\begin{frame}[layout=fullB]{Summary}
  \BTKContent{text}{
    \begin{itemize}
      \item Layouts define named regions.
      \item Regions give normal LaTeX flow inside bounded boxes.
      \item Overlays reference stable anchors for “Power mode”.
    \end{itemize}
  }
\end{frame}
